\documentclass[11pt,a4paper]{article}
\usepackage[margin=1in]{geometry}
\usepackage{amsmath,amssymb,amsfonts}
\usepackage{graphicx}
\usepackage{booktabs}
\usepackage{hyperref}
\usepackage{xcolor}
\usepackage{longtable}
\usepackage{enumitem}
\usepackage[T1]{fontenc}

\hypersetup{colorlinks=true, linkcolor=blue!60!black, urlcolor=blue!50!black, citecolor=green!50!black}

\newcommand{\filepath}[1]{\texttt{\small #1}}
\newcommand{\R}{\mathbb{R}}

\title{\textbf{Paper Explainer Package}\\[0.5em]
\large Comprehensive Technical Companion to:\\[0.3em]
\textit{Gradient-Free Stochastic Optimization of Navier-Stokes Active Contours\\for Unsupervised High-Throughput Phenotyping}\\[0.5em]
\normalsize Chimdi Walter Ndubuisi and Toni Kazic}

\author{Auto-generated companion document --- 2026-03-13}
\date{}

\begin{document}
\maketitle
\tableofcontents
\newpage

%=============================================================================
\section{Overview and Purpose}
%=============================================================================

This document is a \textbf{structured explanatory companion} to the manuscript
``Gradient-Free Stochastic Optimization of Navier-Stokes Active Contours for
Unsupervised High-Throughput Phenotyping'' by Ndubuisi and Kazic.  It is
\textbf{not} a new research paper.  Instead, it serves as:

\begin{itemize}[nosep]
  \item A guided walkthrough of the manuscript from introduction to conclusion
  \item An audit of every experiment output on disk, with file paths
  \item An explanation of the mathematics, methodology, and literature
  \item A figure-by-figure and table-by-table interpretation guide
  \item A provenance record documenting which outputs are final, partial, or diagnostic
\end{itemize}

\paragraph{Manuscript location.}
\filepath{Paper\_3/manuscript\_final\_strengthened/tex/Latest\_Paper\_1\_strengthened.tex}\\
Compiled PDF: \filepath{Paper\_3/manuscript\_final\_strengthened/pdf/Latest\_Paper\_1\_strengthened.pdf}

%=============================================================================
\section{What the Paper Studies}
%=============================================================================

The paper addresses the problem of \textbf{unsupervised segmentation of plant disease lesions} in high-resolution 16-bit leaf images.  Lesions exhibit enormous morphological diversity: sharp necrotic spots, faint chlorotic halos, coalescing regions, and complex backgrounds (venation, glare, dust).

The core insight is to treat the tuning of a \textbf{physics-based segmentation engine} (Navier-Stokes-inspired diffusion + active contour snakes) as a \textbf{per-image black-box optimization problem}, rather than fixing one parameter set for all images.

\subsection{Why This Matters}
\begin{itemize}[nosep]
  \item Deep learning requires thousands of pixel-accurate masks --- unavailable in most phenotyping labs.
  \item Fixed-parameter physics models fail on heterogeneous datasets.
  \item Per-image optimization overcomes both limitations: no labels needed, and parameters adapt to each leaf.
  \item The resulting pseudo-labels can train a fast student network for deployment.
\end{itemize}

%=============================================================================
\section{Core Scientific Objective}
%=============================================================================

Develop a fully automated, unsupervised pipeline that:
\begin{enumerate}[nosep]
  \item Uses PDE-based diffusion to create smooth energy landscapes from image gradients
  \item Evolves active contours (snakes) to segment lesions
  \item Optimizes 7 physics parameters per image via gradient-free stochastic search
  \item Scores segmentation quality without ground truth (heuristic objective)
  \item Produces 16-bit masks and metadata for downstream phenotyping
  \item Distills the optimizer's outputs into a fast UNet student
\end{enumerate}

%=============================================================================
\section{High-Level Method Summary}
%=============================================================================

\begin{enumerate}[nosep]
  \item \textbf{Pre-processing}: 16-bit TIFF $\to$ 8-bit working space; channel fusion (green + grayscale + Lab $b^*$); CLAHE contrast enhancement; leaf masking.
  \item \textbf{Energy field}: Navier-Stokes-inspired diffusion-reaction PDE smooths image gradients into a clean energy map (80 iterations, CFL-stable).
  \item \textbf{Contour evolution}: Threshold energy map $\to$ candidate blobs $\to$ fit snakes with tension $\alpha$, rigidity $\beta$, balloon $\gamma$.
  \item \textbf{Scoring}: Heuristic objective combining gradient alignment, Lab color separability, topology reward, and penalty terms.
  \item \textbf{Optimization}: Random-restart hill climbing (64 random + 48 refine evaluations per image) over bounded 7D parameter space.
  \item \textbf{Output}: Per-image optimized masks, overlays, parameters, and scores.
  \item \textbf{Distillation}: Patchify 174 leaves $\to$ 4,025 training patches $\to$ train UNet student (40 epochs, BCE+Dice loss).
\end{enumerate}

%=============================================================================
\section{Section-by-Section Explanation of the Manuscript}
%=============================================================================

\subsection{Section 1: Introduction}
Motivates the work by identifying the gap between deep learning's data hunger and physics-based models' parameter sensitivity.  Lists four contributions: (1) diffusion engine, (2) no-GT objective, (3) stochastic optimizer, (4) phenotyping-ready pipeline.

\textbf{Note:} Contains a TODO for biological experiment context (genotype, pathogen, DPI, camera) --- this should be filled before submission.

\subsection{Section 2: Related Work}
Five subsections covering: active contours and level sets (Kass, Caselles, Chan-Vese, GVF), PDEs in vision (Navier-Stokes inpainting, PlantCV), modern segmentation (U-Net through SAM), hyperparameter optimization, and plant phenotyping tooling.

The literature positions the paper at the intersection of classical physics-based segmentation and modern automated parameter search, with phenotyping as the application domain.

\subsection{Section 3: Problem Setup and Notation}
Defines the image $I$, parameter vector $\theta = (\mu, \lambda, \alpha, \beta, \gamma, T)$, and the per-image optimization problem $\theta^* = \arg\max S(I;\theta)$.

Introduces two key mathematical concepts:
\begin{itemize}[nosep]
  \item \textbf{Candidate trace sets} $\mathcal{T}_i$ (Eq.~2): logging all evaluated candidates during optimization
  \item \textbf{Projected pairwise surfaces} $\mathcal{S}_{p,q}$ (Eq.~3): projecting the 7D landscape onto 2D parameter planes
  \item \textbf{Parameter-cloud geometry} (Section 3.2): PCA of global/local solutions, centroid shift, $k$-NN analysis
\end{itemize}

\subsection{Section 4: Mathematical Framework}

\subsubsection{Navier-Stokes Diffusion (Eq.~4--5)}
The diffusion-reaction PDE evolves an energy field $E^{(t)}$:
\[
E^{(t+1)} = E^{(t)} + \Delta t \Big( \mu \Delta E^{(t)} + (\lambda+\mu)\,\mathrm{div}(\nabla E^{(t)}) + q \cdot (F - E^{(t)}) \Big)
\]
\begin{itemize}[nosep]
  \item $\mu$: viscosity (Laplacian smoothing)
  \item $(\lambda+\mu)\,\mathrm{div}(\nabla E)$: elastic coupling (coherent structure)
  \item $F = \|(I_x - S_x, I_y - S_y)\|_2$: forcing from seed-vs-image gradient discrepancy
  \item $q = \mathbf{1}[G > T]$: edge-gating mask
  \item Stability: CFL condition $\Delta t \leq h^2/(4\mu)$; with $h=1$, requires $\mu \leq 0.25$
  \item 80 iterations, normalize $E$ to $[0,255]$
\end{itemize}

\subsubsection{Active Contour Dynamics (Eq.~6--7)}
Snake energy functional with internal terms ($\alpha$ tension, $\beta$ rigidity) and external energy from the diffused field.  Balloon force $\gamma$ drives outward expansion.  Multi-lesion handling via energy-map thresholding with scale constraints ($f_{\min} = 5\times10^{-4}$, $f_{\max} = 0.45$).

\subsection{Section 5: Methodology --- Stochastic Parameter Search}

\subsubsection{Heuristic Objective (Eq.~9)}
\[
S(\theta) = w_g \,\mathcal{G}(M,I) + w_c \,\mathcal{C}(M,I) + w_n \log(1+N(M)) - \mathcal{P}_{\mathrm{area}}(M) - \mathcal{P}_{\mathrm{noise}}(M)
\]
\begin{itemize}[nosep]
  \item $\mathcal{G}$: Gradient alignment --- boundaries on stronger-than-average gradients (Eq.~10)
  \item $\mathcal{C}$: Chromatic separability --- Lab $a^*b^*$ distance inside vs outside (Eq.~11)
  \item $\log(1+N)$: Topology reward --- multiple lesions with diminishing returns
  \item $\mathcal{P}_{\mathrm{area}}$: Penalty for empty or too-large masks
  \item $\mathcal{P}_{\mathrm{noise}}$: Penalty for ultra-small dust components
\end{itemize}

\subsubsection{Search Algorithm (Algorithm 1)}
Two-phase: (1) $K=64$ uniform random samples, (2) top-4 seeds refined with $R=12$ perturbation steps each, cooling $\sigma_j = \sigma_0 \cdot 0.85^j$.  Total $\sim$112 evaluations per image.

\subsection{Section 6: Implementation Details}
16-bit handling, channel fusion, CLAHE, scale-aware filtering, downscaling during search ($s=0.75$), timeouts, parallelism.

\subsection{Section 7: Experiments}
Dataset: 174 16-bit TIFF leaf images.  Evaluation: no-GT (all), sparse GT (60 planned), cross-condition.  UNet student: 4,025 patches, 40 epochs.

\subsection{Section 8: Results}
Detailed in Section~\ref{sec:experiments} below.

\subsection{Section 9: Multi-Stage Inference}
Two-stage nested lesion approach: outer (green) + inner (magenta) masks.  Algorithm 2.

\subsection{Section 10: Discussion}
Parameter transfer, morphology fidelity, landscape geometry, manifold structure, limitations.

\subsection{Section 11: Conclusion}
Summarizes contributions and points to amortized prediction as future work.

\subsection{Appendices A--C}
Additional optimizer diagnostics, supplementary trace plots, supplementary geometry plots.

%=============================================================================
\section{Experiment-by-Experiment Explanation}
\label{sec:experiments}
%=============================================================================

\subsection{Experiment 1: Production Optimizer Run (174 Leaves)}

\begin{tabular}{ll}
\toprule
\textbf{What} & Per-image stochastic optimization on all 174 leaves \\
\textbf{Why} & Generate optimized masks and parameters for every image \\
\textbf{How} & Random-restart hill climbing (64 random + 48 refine) \\
\textbf{Output} & \filepath{opt\_summary\_local.csv} \\
\textbf{Key result} & 100\% OK; mean score 21.10; median runtime 242s \\
\textbf{Paper role} & Main paper --- Figs 1--3, referenced throughout \\
\textbf{Status} & FINAL \\
\bottomrule
\end{tabular}

\subsection{Experiment 2: Objective Function Ablations}

\begin{tabular}{ll}
\toprule
\textbf{What} & Remove one objective term at a time \\
\textbf{Why} & Justify each term in the composite objective (Eq.~9) \\
\textbf{How} & Same optimizer, modified weights, 30-image stratified subset \\
\textbf{Output} & \filepath{ablations\_objective/} (6 subdirectories) \\
\textbf{Key result} & Contour-count removal: score 14.80 $\to$ 3.68 (most impactful) \\
 & Noise penalty removal: $n_c$ 134 $\to$ 188 (fragmentation) \\
\textbf{Paper role} & Main paper --- Table I \\
\textbf{Status} & FINAL \\
\bottomrule
\end{tabular}

\subsection{Experiment 3: Search Strategy Ablations}

\begin{tabular}{ll}
\toprule
\textbf{What} & Modify one search knob at a time \\
\textbf{Why} & Justify the two-phase search design \\
\textbf{How} & Same objective, modified search params, 30-image subset \\
\textbf{Output} & \filepath{ablations\_search/} (5 subdirectories) \\
\textbf{Key result} & No-downscale highest score (16.58) but 34\% slower \\
 & Random-only competitive (14.88 vs 14.95) \\
\textbf{Paper role} & Main paper --- Table II \\
\textbf{Status} & FINAL \\
\bottomrule
\end{tabular}

\subsection{Experiment 4: Teacher--Student Distillation}

\begin{tabular}{ll}
\toprule
\textbf{What} & Train UNet student on optimizer pseudo-labels \\
\textbf{Why} & Amortize per-image search into fast inference \\
\textbf{How} & 4,025 patches, 40 epochs, BCE+Dice, threshold 0.30 \\
\textbf{Output} & \filepath{teacher\_student\_metrics.csv}, \\
 & \filepath{teacher\_student\_metrics\_roi\_thr0p30.csv}, \\
 & \filepath{teacher\_student\_figures/}, \filepath{threshold\_sweep.csv} \\
\textbf{Key result} & Baseline Dice 0.598 (all), 0.642 (teacher-nonempty) \\
 & 18 failures, 12 teacher-empty disagreements \\
\textbf{Paper role} & Main paper --- Section 8.1, Figs 4--6 \\
\textbf{Status} & FINAL \\
\bottomrule
\end{tabular}

\subsection{Experiment 5: 9-Leaf Representative Trace Study}

\begin{tabular}{ll}
\toprule
\textbf{What} & Instrumented rerun with candidate-level trace logging \\
\textbf{Why} & Characterize objective landscape geometry \\
\textbf{How} & 9 selected leaves (3 high/3 median/3 low score); \\
 & identical settings as production; SLURM job 12783426 \\
\textbf{Output} & \filepath{objective\_surface\_traces\_representative\_fullbudget\_9leaf/} \\
\textbf{Provenance} & ALL traces from single dedicated rerun --- NO fallback \\
\textbf{Key result} & High-score: sharp landscapes, refine gain +5.77 \\
 & Low-score: flat landscapes, refine gain +0.56 \\
 & Dominant pair: $T \times \beta$ (7/9 cases) \\
\textbf{Paper role} & Main paper --- Figs 9--12, Section 8.4 \\
\textbf{Status} & FINAL --- homogeneous dedicated rerun \\
\bottomrule
\end{tabular}

\subsection{Experiment 6: Parameter-Cloud Manifold Geometry}

\begin{tabular}{ll}
\toprule
\textbf{What} & PCA/t-SNE/UMAP analysis of 7D parameter space \\
\textbf{Why} & Characterize structure of optimizer's solution space \\
\textbf{How} & 150 global/local parameter pairs from 15 leaves \\
\textbf{Output} & \filepath{parameter\_shift\_manifold\_geometry/} \\
\textbf{Key result} & 2 PCs capture 46.5\% variance; $k$-NN 0.534 (vs 0.33 random) \\
\textbf{Paper role} & Main paper --- Figs 13--14, Section 8.5 \\
\textbf{Status} & FINAL \\
\bottomrule
\end{tabular}

\subsection{Experiment 7: Solid-3D Parameter Geometry}

\begin{tabular}{ll}
\toprule
\textbf{What} & Covariance ellipsoids, convex hulls, density isosurfaces \\
\textbf{Why} & Volumetric characterization of parameter clouds \\
\textbf{How} & 2-sigma ellipsoids in PCA space; hull analysis \\
\textbf{Output} & \filepath{parameter\_shift\_3d\_analysis\_solid/} \\
\textbf{Key result} & Centroid shift 1.616 std units; partial $\Delta D$ separation \\
\textbf{Paper role} & Main paper --- Fig 15; Appendix C \\
\textbf{Status} & FINAL \\
\bottomrule
\end{tabular}

\subsection{Experiment 8: Patch Morphology Transfer}

\begin{tabular}{ll}
\toprule
\textbf{What} & Test whether per-leaf params transfer to 512$\times$512 patches \\
\textbf{Why} & Validate global-to-local parameter stability \\
\textbf{How} & 150 patches from 15 leaves; global vs local re-optimization \\
\textbf{Output} & \filepath{patch\_local\_opt\_v2.csv}, \\
 & \filepath{patch\_transfer\_morphology\_eval.csv}, \\
 & \filepath{patch\_transfer\_morphology\_plots/} \\
\textbf{Key result} & Median Dice 0.78; area underestimated (0.026 vs 0.169); \\
 & local re-opt reduces count error 35\% \\
\textbf{Paper role} & Main paper --- Table III, Figs 17--19, Discussion \\
\textbf{Status} & FINAL \\
\bottomrule
\end{tabular}

%=============================================================================
\section{Figure-by-Figure Explanation with File Paths}
%=============================================================================

\subsection{Main Text Figures}

\subsubsection{Fig 1 --- Optimizer Status Bar}
\filepath{optimizer\_summary/status\_bar.png}\\
Shows bar chart of optimizer run status for all 174 images.  All bars show ``OK'' status.\\
\textbf{Takeaway}: 100\% convergence proves pipeline reliability across diverse phenotypes.

\subsubsection{Fig 2 --- Runtime Histogram}
\filepath{optimizer\_summary/elapsed\_seconds\_hist.png}\\
Histogram of per-image wall-clock time (seconds).  Median $\sim$250s.\\
\textbf{Takeaway}: Per-image optimization takes $\sim$4 minutes on CPU --- practical for 174 images.

\subsubsection{Fig 3 --- Score Histogram}
\filepath{optimizer\_summary/score\_hist.png}\\
Distribution of objective scores across 174 images.  Range 11--46, mean 21.10.\\
\textbf{Takeaway}: Scores are unimodal with no catastrophic failures.

\subsubsection{Fig 4 --- Threshold Sweep}
\filepath{threshold\_sweep\_plot.png}\\
Dice as a function of probability threshold on stitched UNet outputs.\\
\textbf{Takeaway}: Best at 0.3; teacher-empty disagreements persist regardless of threshold.

\subsubsection{Figs 5--6 --- Teacher-Student Montages}
\filepath{teacher\_student\_figures/montage\_teacher\_student.png}\\
\filepath{paper\_disagreement\_montage\_roi\_thr0p30.png}\\
Fig 5: worst/median/best comparisons.  Fig 6: teacher-empty disagreement cases.\\
\textbf{Takeaway}: Student approximates teacher well; failures are semantic (inner-mask definition), not background leakage.

\subsubsection{Figs 7--8 --- Parameter-Morphology Correlations}
\filepath{optimizer\_summary/ (derived scatter plots)}\\
Fig 7: Energy threshold vs.\ lesion count (inverse).  Fig 8: Balloon force vs.\ area fraction (positive).\\
\textbf{Takeaway}: Optimized parameters have physically interpretable relationships with output morphology.

\subsubsection{Figs 9--11 --- Trace-Based Objective Surfaces}
\filepath{objective\_surface\_traces\_representative\_fullbudget\_9leaf/surface\_plots/}\\
Projected $T$--$\beta$ surfaces for high (DSC\_0198), median (DSC\_0379), and low (DSC\_0261) score leaves.\\
\textbf{Takeaway}: Landscape curvature decreases with phenotypic difficulty.  High-score leaves have sharp optima; low-score leaves are flat.

\subsubsection{Fig 12 --- Convergence Trajectory}
\filepath{objective\_surface\_traces\_representative\_fullbudget\_9leaf/trajectory\_plots/}\\
Best-score-so-far vs.\ candidate index for DSC\_0198 (high-score).\\
Blue = random phase, orange = refine phase.\\
\textbf{Takeaway}: Clear phase transition; refinement adds measurable gain on favorable landscapes.

\subsubsection{Fig 13 --- PCA 2D Global-Local Projection}
\filepath{parameter\_shift\_manifold\_geometry/plots/pca2d\_global\_local\_dD.png}\\
Global (circles) and local (triangles) parameters in PCA space, colored by $\Delta D$.\\
\textbf{Takeaway}: PC1 separates the two clouds; centroid shift magnitude 1.616 standardized units.

\subsubsection{Fig 14 --- Mean Parameter Shift Profile}
\filepath{parameter\_shift\_manifold\_geometry/plots/mean\_delta\_parameters.png}\\
Bar chart of mean shift per parameter.\\
\textbf{Takeaway}: Energy threshold dominates --- local re-optimization primarily adjusts edge sensitivity.

\subsubsection{Fig 15 --- Covariance Ellipsoids}
\filepath{parameter\_shift\_3d\_analysis\_solid/pca3d\_covariance\_ellipsoids.png}\\
2-sigma ellipsoids for global (blue) and local (red) clouds in PCA space.\\
\textbf{Takeaway}: Overlapping but displaced ellipsoids confirm structured parameter organization.

\subsubsection{Fig 16 --- Qualitative Grid}
\filepath{teacher\_student\_figures/ (qual\_grid)}\\
6-row grid from Dice$\approx$0 to Dice$\approx$0.84.  Columns: input, teacher overlay, student overlay, masks.\\
\textbf{Takeaway}: Visual evidence spanning the full quality spectrum.

\subsubsection{Figs 17--19 --- Morphology Transfer Plots}
\filepath{patch\_transfer\_morphology\_plots/hist\_area\_frac\_error.png}\\
\filepath{patch\_transfer\_morphology\_plots/hist\_count\_error.png}\\
\filepath{patch\_transfer\_morphology\_plots/scatter\_teacher\_vs\_pred\_area\_frac.png}\\
Area error histogram, count error histogram, teacher vs.\ predicted scatter.\\
\textbf{Takeaway}: Both parameter sets underestimate area; local re-opt reduces count error 35\%.

\subsection{Appendix Figures}

\subsubsection{Appendix A --- Optimizer Output Diagnostics}
\filepath{optimizer\_summary/n\_contours\_hist.png} and \filepath{optimizer\_summary/area\_px\_hist.png}\\
Connected-component count and mask area distributions.  Seven parameter histograms (one per parameter).

\subsubsection{Appendix B --- Additional Trace Trajectories}
\filepath{objective\_surface\_traces\_representative\_fullbudget\_9leaf/trajectory\_plots/}\\
Six additional convergence trajectories (2 per tier: DSC\_0166+0163, DSC\_0274+0205, DSC\_0185+0199).

\subsubsection{Appendix C --- Supplementary Geometry}
\filepath{parameter\_shift\_3d\_analysis\_solid/pca3d\_convex\_hull\_global\_local.png}\\
\filepath{parameter\_shift\_3d\_analysis\_solid/pca3d\_positive\_vs\_negative\_dD\_surfaces.png}\\
\filepath{parameter\_shift\_manifold\_geometry/plots/tsne2d\_delta\_ddgroup.png}\\
\filepath{parameter\_shift\_manifold\_geometry/plots/pca\_centroid\_shift\_groups.png}\\
Convex hulls, $\Delta D$ surfaces, t-SNE, centroid shift by group.

%=============================================================================
\section{Table-by-Table Explanation with File Paths}
%=============================================================================

\subsection{Table I --- Objective Ablations}
\textbf{Source}: \filepath{ablations\_objective/} (6 subdirectories)\\
\textbf{Section}: Results 8.3\\
\textbf{Format}: 6 rows (full + 5 ablated variants) $\times$ 6 columns (OK\%, Non-empty\%, Med.\ score, Med.\ area, Med.\ $n_c$, Med.\ time)

\textbf{How to read}: Compare each ablated row against the ``Full objective (ours)'' row.  Arrows indicate preferred direction ($\uparrow$ or $\downarrow$).

\textbf{Key findings}:
\begin{itemize}[nosep]
  \item Removing contour count ($w_{\mathrm{topo}}=0$): score drops from 14.80 to 3.68 --- \textbf{most impactful term}
  \item Removing noise penalty ($w_{\mathrm{small}}=0$): $n_c$ jumps from 134 to 188 --- fragmentation
  \item Removing gradient alignment ($w_{\mathrm{grad}}=0$): score drops to 9.71
  \item Removing color distance or area penalty: modest effects
\end{itemize}

\subsection{Table II --- Search Strategy Ablations}
\textbf{Source}: \filepath{ablations\_search/} (5 subdirectories)\\
\textbf{Section}: Results 8.3\\
\textbf{Format}: 5 rows $\times$ 6 columns (same as Table I)

\textbf{Key findings}:
\begin{itemize}[nosep]
  \item No-downscale: highest score (16.58) but 34\% more runtime (5427.9s vs 4046.5s)
  \item Random-only: surprisingly competitive (14.88 vs 14.95)
  \item Low budget ($K=16$): still achieves 14.91 --- diminishing returns from more random samples
  \item Fast timeout (60s): modest score reduction (14.44)
\end{itemize}

\subsection{Table III --- Patch Morphology Transfer}
\textbf{Source}: \filepath{patch\_transfer\_morphology\_eval.csv} + \filepath{patch\_local\_opt\_v2.csv}\\
\textbf{Section}: Discussion\\
\textbf{Format}: Metrics for All/Interior/Boundary patches, comparing global vs.\ local parameters against teacher

\textbf{Key findings}:
\begin{itemize}[nosep]
  \item Dice is modest: global 0.430, local 0.443 ($+$0.013)
  \item Area fraction severely underestimated: 0.026 vs.\ teacher 0.169
  \item Count error reduced 35\% by local re-optimization: 19.1 $\to$ 12.4
  \item Interior patches benefit more from local re-opt than boundary patches
  \item \textbf{Conclusion}: Dice alone is insufficient for phenotyping quality; morphological metrics matter
\end{itemize}

%=============================================================================
\section{Literature and Mathematics Explained}
%=============================================================================

\subsection{Literature Landscape}

The paper draws on five distinct literature threads:

\begin{enumerate}[nosep]
  \item \textbf{Active contours}: Kass et al.\ (1988) $\to$ Caselles (1997, geodesic) $\to$ Chan-Vese (2001, region-based) $\to$ GVF (Xu \& Prince, 1998).  The paper extends this tradition by adding automated per-image parameter search.

  \item \textbf{PDEs in vision}: Bertalmio et al.\ (2001, Navier-Stokes inpainting) inspires the diffusion formulation.  The paper's energy-field diffusion is a scalar version of GVF with fluid-dynamics-inspired elastic coupling.

  \item \textbf{Modern segmentation}: U-Net (2015) through SAM (2023).  Used as context and baselines.  The student model is U-Net-based.

  \item \textbf{Black-box optimization}: Yu \& Zhu (2020) survey; Hansen \& Ostermeier (2001, CMA-ES); Jones et al.\ (1998, EGO).  Justifies gradient-free approach for non-differentiable pipelines.

  \item \textbf{Plant phenotyping}: Gehan \& Kellogg (2017), Mahlein (2016), PlantCV (2025), Singh et al.\ (2016).  Establishes the application domain and emphasizes trait-producing pipelines.
\end{enumerate}

\subsection{Mathematical Components}

\textbf{11 numbered equations} in the manuscript:
\begin{itemize}[nosep]
  \item Eq.~1, 8: Per-image optimization formulation
  \item Eq.~2: Candidate trace set definition
  \item Eq.~3: Projected pairwise surface definition
  \item Eq.~4: Diffusion-reaction PDE update
  \item Eq.~5: Forcing term from gradient discrepancy
  \item Eq.~6: Snake energy functional
  \item Eq.~7: Snake evolution equation
  \item Eq.~9: Composite heuristic objective
  \item Eq.~10: Gradient alignment term
  \item Eq.~11: Chromatic separability term
\end{itemize}

Plus the necrosis ratio formula in Section 9 and the Dice/overlap formulations in the evaluation.

%=============================================================================
\section{Main Findings and What They Mean}
%=============================================================================

\begin{enumerate}
  \item \textbf{100\% convergence across 174 diverse leaves}: The pipeline is robust and does not fail catastrophically on any image.

  \item \textbf{Ablations justify every component}: Contour-count topology reward is the most impactful objective term; the full search pipeline balances quality and speed.

  \item \textbf{Teacher-student distillation works} (median Dice 0.682 on teacher-nonempty images), but 12 teacher-empty disagreements indicate a semantic gap between inner-mask and outer-lesion definitions.

  \item \textbf{Score-tier-dependent landscape geometry}: The objective landscape is sharp for morphologically favorable images (high scores) and flat for challenging ones (low scores).  This has practical implications for adaptive budget allocation.

  \item \textbf{Structured parameter manifold}: 46.5\% of variance in 2 PCs; shift direction partially predicts optimization success ($k$-NN 0.534 vs.\ 0.33 random).  This supports amortized parameter prediction.

  \item \textbf{Morphological fidelity diverges from overlap metrics}: Dice alone underestimates issues; area fraction is severely underestimated under patchification (0.026 vs.\ 0.169).  Local re-optimization helps with count (35\% error reduction) but not area.
\end{enumerate}

%=============================================================================
\section{Strongest Evidence}
%=============================================================================

\begin{enumerate}[nosep]
  \item \textbf{Production run} (\filepath{opt\_summary\_local.csv}): 100\% convergence on 174 heterogeneous leaves
  \item \textbf{Ablation tables} (Tables I--II): Each component rigorously justified on 30-image subset
  \item \textbf{9-leaf trace surfaces} (Figs 9--12): Clean provenance (single SLURM job); score-tier pattern
  \item \textbf{Manifold geometry} (Figs 13--15): Quantified structure with $k$-NN above chance
  \item \textbf{Morphology transfer} (Table III): Reveals Dice-morphology divergence important for phenotyping
\end{enumerate}

%=============================================================================
\section{What is Provisional, Appendix-Only, or Diagnostic}
%=============================================================================

\subsection{Provisional / Incomplete}
\begin{itemize}[nosep]
  \item \filepath{objective\_surface\_traces\_representative\_fullbudget\_plus/}: manifold and solid-3D subdirs EMPTY
  \item \filepath{param\_predictor/}: EMPTY --- future work (amortized prediction)
  \item Biological experiment context TODO in manuscript (genotype, pathogen, DPI)
  \item Sparse GT evaluation on 60-image subset: protocol described but annotation not completed
  \item Baseline comparisons (Otsu, labeled U-Net, SAM): recommended but not reported
\end{itemize}

\subsection{Appendix-Only}
\begin{itemize}[nosep]
  \item Parameter histograms (Fig A2): useful for reproducibility, not primary evidence
  \item Additional convergence trajectories (Appendix B): extend Fig 12 pattern
  \item Supplementary geometry (Appendix C): t-SNE, hulls, centroid groups
  \item Full-population traces (\filepath{objective\_surface\_traces\_fullstudy/}): population evidence
\end{itemize}

\subsection{Diagnostic Only --- Not for Paper}
\begin{itemize}[nosep]
  \item Codex-style variants (\filepath{optimization\_surface\_analysis\_codex/}, \filepath{parameter\_shift\_3d\_analysis\_solid\_codex/})
  \item Per-leaf montage PNGs (\filepath{montage\_patch\_opt\_v2\_*.png})
  \item \filepath{postprocess\_sweep.csv}: parameter sweep diagnostic
  \item \filepath{optimizer\_full/}, \filepath{optimizer\_full\_alt/}: empty directories
  \item \filepath{experiment\_summary\_master/}: prior summary attempt
\end{itemize}

%=============================================================================
\section{Manifold / Solid-3D / Optimization Geometry Assessment}
%=============================================================================

\textbf{Question}: Are manifold and solid-3D geometry figures already present?

\textbf{Answer}: \textbf{Yes} --- all geometry figures are already present in the outputs and are included in the current manuscript.

\begin{itemize}[nosep]
  \item \textbf{PCA manifold figures}: Present in \filepath{parameter\_shift\_manifold\_geometry/plots/} (30 files) and copied to \filepath{figs\_from\_outputs/manifold/} (6 files used in paper)
  \item \textbf{Solid-3D figures}: Present in \filepath{parameter\_shift\_3d\_analysis\_solid/} (25 files) and copied to \filepath{figs\_from\_outputs/solid3d/} (4 files used in paper)
  \item \textbf{Merged final figures}: Present in \filepath{optimization\_manifold\_and\_solid3d\_final/figures/} (19 files)
  \item \textbf{9-leaf trace geometry}: Present in \filepath{objective\_surface\_traces\_representative\_fullbudget\_9leaf/} and copied to \filepath{figs\_from\_outputs/9leaf/} (54 files)
\end{itemize}

\textbf{No new geometry figures needed.}  All required manifold, PCA, UMAP, t-SNE, solid-3D hull, ellipsoid, and density plots already exist on disk and are reflected in the manuscript.

%=============================================================================
\section{Directory Structure Summary}
%=============================================================================

Key paths (all relative to \filepath{experiments/exp\_v1/outputs/}):

\begin{verbatim}
experiments/exp_v1/outputs/
  opt_summary_local.csv                          <- CORE: 174-leaf results
  ablations_objective/                           <- Table I source
  ablations_search/                              <- Table II source
  objective_surface_traces_representative_fullbudget_9leaf/
    candidate_traces/                            <- 9 trace CSVs
    surface_plots/                               <- 27+ surface plots
    trajectory_plots/                            <- 54+ trajectory plots
    reports/                                     <- surface_report.pdf/.txt
  parameter_shift_manifold_geometry/
    plots/                                       <- 30 manifold plots
    reports/                                     <- 3 reports
    data/                                        <- 4 CSVs
  parameter_shift_3d_analysis_solid/             <- 25 solid-3D files
  optimization_manifold_and_solid3d_final/       <- 19 merged figures
  optimizer_summary/                             <- 17 summary files
  teacher_student_figures/                       <- montages + plots
  patch_transfer_morphology_plots/               <- 7 morph plots
  teacher_student_metrics.csv                    <- baseline metrics
  teacher_student_metrics_roi_thr0p30.csv        <- ROI-gated metrics
  threshold_sweep.csv + threshold_sweep_plot.png <- sweep data + plot
  failures_dice_lt_0p3.csv                       <- 18 failure cases
  teacher_empty_cases.csv                        <- 12 disagreements
  patch_local_opt_v2.csv                         <- 150 patch results
  patch_transfer_morphology_eval.csv             <- morphology metrics
  patches/                                       <- training data
\end{verbatim}

%=============================================================================
\section{Final Summary}
%=============================================================================

This paper presents a novel approach to unsupervised lesion segmentation that bridges physics-based modeling and automated parameter optimization.  The key innovation is treating per-image physics parameter tuning as a black-box optimization problem, enabling robust segmentation across 174 morphologically diverse leaf images without any ground-truth labels.

The experimental evidence is comprehensive:
\begin{itemize}[nosep]
  \item Production run achieves 100\% convergence (strong reliability)
  \item Ablations justify every objective and search component
  \item Trace-based landscape analysis reveals score-tier-dependent geometry
  \item Manifold analysis confirms structured, low-dimensional parameter organization
  \item Teacher-student distillation enables fast inference (with known semantic limitations)
  \item Morphology transfer analysis reveals important Dice-morphology divergence
\end{itemize}

The main gaps are: (1) no quantitative baseline comparisons, (2) sparse GT annotation not yet completed, (3) biological context TODO unfilled, (4) amortized parameter predictor not yet built.

All experiment outputs have clean provenance.  The 9-leaf trace study (the most novel analysis) uses a single homogeneous SLURM rerun with no fallback or reused traces.  No partial or mixed-provenance results are presented as final in the main paper.

\end{document}
